% AC 最优潮流：Native IPM 求解器
% 本章文件由 \input 引入，不含导言区。
\section{AC 最优潮流：Native IPM 求解器}\label{sec:acopfnative}

本章给出平台自研 AC 最优潮流（AC OPF）引擎的技术口径。求解入口为
\allowbreak\srcpath{hacdcpf::opf::solve\_ac\_opf}\allowbreak （声明于
\srcpath{include/hacdcpf/optimal\_power\_flow/ac\_opf\_solver.hpp}:42--43，
实现于 \srcpath{src/optimal\_power\_flow/ac\_opf.cpp}:3038--4209）。该实现包含两条
自研数值路径：一条是内嵌于 \srcpath{ac\_opf.cpp} 的紧致变量原生内点循环
（结果标记为 \fld{OPFSolverPath::\allowbreak NativeAC}），另一条是作用于 Parity 全空间
建模的原始对偶内点核 \srcpath{parity::solve\_primal\_dual\_ipm}\allowbreak （接口声明于
\srcpath{include/hacdcpf/optimal\_power\_flow/native\_ipm\_solver.hpp}:90，
结果标记为 \fld{OPFSolverPath::\allowbreak ParityIPM}）；两者共享同一组选项与结果类型，
并可在编译入 Ipopt 时自动回退到嵌入式 Ipopt。需要特别说明：尽管接口文档以
``polar AC balance'' 描述核心等式（\srcpath{ac\_opf\_solver.hpp}:18--31），
两条路径的电压变量均为\textbf{极坐标}（幅值 $V$ 与相角 $\theta$），而非直角坐标
实虚部；本章全部公式以代码实际采用的极坐标形式给出。

\subsection{功能定位与求解路径}\label{subsec:acopfnative:role}

\subsubsection*{功能定位}
Native IPM 是平台交直流混合 OPF 的自研引擎层，承担以下职责
（\srcpath{src/optimal\_power\_flow/ac\_opf.cpp}:3038--4209）：
\begin{itemize}
  \item 对 rich 模型 \fld{HybridPowerSystem} 做合法性门禁（空 AC 母线、LCC 就绪性、
        参考母线资格、图拓扑孤岛），随后经 \srcpath{core::make\_solver\_data}
        装配 canonical \fld{SolverData}（含 Ybus 与直流电导矩阵 \fld{gdc}）；
  \item 在标幺体系（$S_{\mathrm{base}}$ 取算例 \fld{base\_mva}，典型 100~MVA）下组装
        非线性规划：极坐标节点功率平衡、直流节点平衡、换流器耦合等式，
        支路视在功率与换流器容量圆不等式，以及全体决策变量的盒式界；
  \item 以原始对偶内点框架求解，输出电压、发电调度、换流器功率、LMP 及
        诚实口径的 \fld{model\_limitations} / \fld{converter\_model\_scope} 声明。
\end{itemize}
与 Parity/Ipopt 的关系：Parity 全空间建模（见第~\ref{sec:parityform}~章）是
\textbf{生产推荐的建模层}，覆盖 VSC 控制模式、直流电压等式、DER 全家族与 LCC
准稳态；其内层牛顿引擎默认就是本章的 Native IPM 核
\srcpath{parity::solve\_primal\_dual\_ipm}（见第~\ref{sec:parityipm}~章），
Ipopt 仅作为可选回退。代码自身把内嵌于 \srcpath{ac\_opf.cpp} 的紧致路径标记为
legacy：结果中会写入建议改用 ParityIPM 的 \fld{model\_limitations} 条目
（\srcpath{src/optimal\_power\_flow/ac\_opf.cpp}:3219--3222）。Ipopt 的可用性由
编译宏 \srcpath{HACDCPF\_HAVE\_IPOPT} 门控（同文件:33--35、2267--2373），
并可用环境变量 \srcpath{HACDCPF\_DISABLE\_IPOPT\_OPF} 在运行时强制关闭
（同文件:2391--2396）。

\subsubsection*{求解路径选择}
\fld{ACOPFOptions::ac\_solver\_backend}（枚举 \fld{ACOPFSolverBackend}，
\srcpath{include/hacdcpf/optimal\_power\_flow/opf\_options.hpp}:28--33）给出四种
后端：\fld{Auto}、\fld{ParityIPM}、\fld{Ipopt}、\fld{EconomicDispatch}。入口处的
路径翻译规则为（\srcpath{src/optimal\_power\_flow/ac\_opf.cpp}:3309--3350）：
\begin{itemize}
  \item \fld{ParityIPM}：置 \fld{enable\_primal\_dual} 与 \fld{use\_parity\_ipm}
        两个标志；内层固定为自研核 \fld{ParityInnerSolver::\allowbreak NativeIPM}，
        不经 Ipopt 翻译；
  \item \fld{Ipopt}：同样进入 Parity 建模，内层固定为 \fld{ParityInnerSolver::Ipopt}，
        未编译 Ipopt 时回退自研核；
  \item \fld{EconomicDispatch}：清除 primal-dual 标志，强制纯电交流
        ``经济调度 + 单次 AC 潮流'' 快速路径；
  \item \fld{Auto}：保持调用方标志；只要算例含任何 DC/VSC/LCC 组件
        （判定见 \srcpath{contains\_hybrid\_acdc\_components}，同文件:3022--3034）
        或存在在运 LCC，即自动升级为 Parity 路径（同文件:3336--3350）。
        纯 AC 且未显式开启 primal-dual 时走快速路径（同文件:3501--3526）；
        纯 AC 且 \fld{enable\_primal\_dual=true}、\fld{use\_parity\_ipm=false}
        时才进入本章重点描述的 legacy 原生内点循环（同文件:3528--4098）。
\end{itemize}
整体决策流程如图~\ref{fig:acopfnative:flow}~所示。

\begin{figure}[H]
\centering
\begin{tikzpicture}[
  font=\footnotesize,
  node distance=0.55cm and 0.9cm,
  box/.style={draw, rectangle, align=center, minimum height=0.85cm, inner sep=3pt},
  dec/.style={draw, diamond, aspect=2.2, align=center, inner sep=1.5pt},
  arr/.style={-{Stealth[length=2.5mm]}, thick},
  lbl/.style={font=\scriptsize, midway, above}
]
\node[box] (entry) {入口校验\\空 AC / LCC 就绪 / 参考母线 / 拓扑};
\node[dec, below=of entry] (hybrid) {含 DC/VSC/LCC？\\或显式 Parity/Ipopt};
\node[box, right=2.6cm of hybrid] (legacy) {legacy 原生障碍--罚\\内点循环（NativeAC）};
\node[box, below=of hybrid] (parity) {Parity 全空间建模\\内层：Native IPM / Ipopt};
\node[box, left=2.6cm of hybrid] (ed) {经济调度 + AC PF\\（EconomicDispatch）};
\node[box, below=of parity] (homo) {目标同伦延续\\（\fld{objective\_homotopy}，可选）};
\node[box, below=of legacy] (fb) {ED+PF 回退\\（仅纯 AC）};
\node[box, below=of homo] (outp) {结果反投影 / 外部电网拆分 /\\ \fld{model\_limitations} 汇总};
\draw[arr] (entry) -- (hybrid);
\draw[arr] (hybrid) -- node[lbl]{否，且未开 primal-dual} (ed);
\draw[arr] (hybrid) -- node[lbl]{否，开 primal-dual 非 parity} (legacy);
\draw[arr] (hybrid) -- node[lbl, right=1pt]{是} (parity);
\draw[arr] (parity) -- (homo);
\draw[arr] (legacy) -- node[lbl, right=1pt]{未收敛} (fb);
\draw[arr] (homo) -- (outp);
\draw[arr] (ed.south) |- (outp.west);
\draw[arr] (fb.south) |- (outp.west);
\end{tikzpicture}
\caption{\srcpath{solve\_ac\_opf} 的求解路径决策流程（依据
\srcpath{src/optimal\_power\_flow/ac\_opf.cpp}:3309--3526 整理）}
\label{fig:acopfnative:flow}
\end{figure}

\subsubsection*{入口门禁}
在进入任何数值路径之前，求解器执行四组静态检查并\textbf{直接拒绝}不合格算例：
\begin{enumerate}
  \item 空 AC 母线集：返回 \fld{status}=``AC OPF failed: empty AC bus set.''
        （\srcpath{src/optimal\_power\_flow/ac\_opf.cpp}:3089--3092）；
  \item LCC 就绪性：对在运 LCC 逐一校验控制模式配对、\fld{n\_bridges} 与
        \fld{vn\_ac\_kv} 正值性、恒 $\alpha$/恒 $\gamma$ 控制的 \fld{x\_comm\_ohm}
        正值性、交流/直流端母线可解析、直流基准电压可得；任一失败即拒绝并写入
        \fld{infeasibility\_hints}（同文件:3094--3140）。注释明确说明动机：潮流的
        LCC 注入助手对不可用站点返回零以便诊断，而 OPF 静默丢弃固定传输量会改变
        可行性与最优解，故必须前置拒绝（同文件:3094--3098）；
  \item 参考母线资格：\srcpath{validation::validate\_reference\_bus\_eligibility}
        存在错误或``裸理想 DC\_V 边界''警告时拒绝，后者附加说明理想 DC\_V 边界
        没有调度成本模型，优化需要显式可控的直流平衡设备（同文件:3142--3161）；
  \item 图拓扑预检：\srcpath{graph::build\_power\_system\_graph} +
        \srcpath{analyze\_topology}，无任何含松弛母线的有效孤岛时判定不可行；
        对无松弛孤岛写入不可服务提示（同文件:3280--3307）。
\end{enumerate}

\subsubsection*{外部电网提升}
每个在运 \fld{ExternalGrid} 在求解前被提升为一台合成 \fld{Generator} 追加到
工作副本：\fld{vg\_pu} 取外部电网 \fld{vm\_pu}，P/Q 上下界取
$\pm\max(1.0,\;$\fld{s\_sc\_max\_mva}$,\;$\fld{base\_mva}$)$~MW/MVAr，成本系数
\fld{cost\_c0/c1/c2} 原样拷贝（\srcpath{src/optimal\_power\_flow/ac\_opf.cpp}:
3163--3196）。代码注释如实声明：短路容量并非进口限值，只是一个``有用的有限
尺度''，以防止过大的合成界使小算例的 KKT 系统病态（同文件:3173--3177）。
求解完成后由 \srcpath{separate\_external\_grid\_dispatch} 把合成机结果拆回
\fld{external\_grid\_p\_mw}/\fld{external\_grid\_q\_mvar}，并把
\fld{pg\_mw}/\fld{qg\_mvar}/\fld{gen\_map} 还原为仅含用户自建发电机的顺序
（同文件:3203--3259）。

\subsection{数学模型}\label{subsec:acopfnative:model}

以下数学模型为 legacy 紧致路径（\fld{NativeAC}）实际求解的形式；Parity 全空间
模型的扩展变量族（DER、储能、柔性负荷、能量路由器、LCC 等）见
第~\ref{sec:parityform}~章。除注明外，功率与电压均为标幺值
（$S_{\mathrm{base}}=$~\fld{base\_mva}）。

\subsubsection*{决策变量与索引布局}
变量向量 $x\in\mathbb{R}^{n}$ 由 \srcpath{build\_index} 按固定块序布局
（\srcpath{src/optimal\_power\_flow/ac\_opf.cpp}:225--344，结构体
\srcpath{ACOPFIndex}:93--142）：
\begin{align*}
x = \left[\,
\underbrace{\theta}_{n_p},\;
\underbrace{V}_{n_q},\;
\underbrace{P_g}_{n_g},\;
\underbrace{Q_g}_{n_g},\;
\underbrace{V_{dc}}_{n_{dc}},\;
\underbrace{P_{ac}}_{n_c},\;
\underbrace{Q_{ac}}_{n_c},\;
\underbrace{P_{dc}}_{n_c}
\,\right]
\end{align*}
各块偏移由 \fld{i\_theta}=0、\fld{i\_vm}=$n_p$ 起顺序累加（同文件:273--274、
326--332），总维数
\begin{align*}
n = n_p + n_q + 2n_g + n_{dc} + 3n_c
\end{align*}
（同文件:332）。其中 $n_p$ 为非松弛交流母线数、$n_q$ 为 PQ 母线数、$n_g$ 为
在运发电机数、$n_{dc}$ 为直流母线数、$n_c$ 为在运 VSC 换流器数。母线归属规则：
松弛母线不设 $\theta$ 变量；仅 PQ 母线设 $V$ 变量（同文件:261--269）；未找到
显式 SLACK 时退化取第 0 号母线为松弛（\srcpath{find\_slack\_bus}，
同文件:178--185）。发电机与换流器仅在 \fld{in\_service} 且端母线可解析时进入
变量集（同文件:287--318；组件 \fld{bus} 字段为 1-based，代码内统一减 1）。
PV/松弛母线的电压幅值、松弛母线的相角不作为变量，而是由参考向量固定：
\fld{vm\_reference\_by\_bus} 先取母线 \fld{vm\_pu}，再被该机端母线发电机的
\fld{vg\_pu} 覆盖（同文件:243--256），解包时先铺参考值再覆盖变量块
（\srcpath{unpack\_state}，同文件:466--538）。

等式右端向量 $g(x)\in\mathbb{R}^{m}$ 的块布局为
（同文件:334--338）：
\begin{align*}
g = \left[\,
\underbrace{g^{P}}_{n_p},\;
\underbrace{g^{Q}}_{n_q},\;
\underbrace{g^{dc}}_{n_{dc}},\;
\underbrace{g^{conv}}_{n_c}
\,\right],\qquad
m = n_p + n_q + n_{dc} + n_c
\end{align*}
故优化自由度恒为 $n-m = 2n_g + 2n_c$：发电机 P/Q 与换流器 P/Q 是仅有的
真正自由度，$V_{dc}$ 与直流平衡一一对应。

\subsubsection*{变量界与初值}
初值与盒式界由 \srcpath{initialize\_bounds\_and\_state} 给出
（\srcpath{src/optimal\_power\_flow/ac\_opf.cpp}:366--464），所有界先经
\srcpath{sanitize\_bounds} 清洗：非有限值回退默认、上下颠倒交换、宽度不足
\fld{kMinInteriorWidth}=$10^{-6}$ 时以中点撑开（同文件:187--203，常数定义于
:41--44）。分块口径：
\begin{itemize}
  \item 相角 $\theta\in[-\pi,\pi]$（硬编码），初值取母线 \fld{va\_deg} 化弧度
        （同文件:375--382）；
  \item 电压幅值 $V\in[$\fld{vmin\_pu}$,$\fld{vmax\_pu}$]$，缺省回退 $[0.5,1.5]$，
        初值取 \fld{vm\_pu}（同文件:384--392）；
  \item 发电机 $P_g\in[$\fld{pmin\_mw}$/S_{\mathrm{base}},\,$
        \fld{pmax\_mw}$/S_{\mathrm{base}}]$，$Q_g$ 同理取
        \fld{qmin\_mvar}/\fld{qmax\_mvar}，缺省回退 $\pm$\fld{kHugeBound}
        （$10^4$），初值取 \fld{pg\_mw}/\fld{qg\_mvar}（同文件:394--413）；
  \item 直流电压普通母线 $V_{dc}\in[0.8,1.2]$（硬编码），DC\_V 型母线收缩为
        设定值 $\pm0.05$，初值取 \fld{vm\_pu}（同文件:415--428）；
  \item 换流器 $P_{ac}$、$Q_{ac}$ 取自身 \fld{pmin\_mw}/\fld{pmax\_mw}、
        \fld{qmin\_mvar}/\fld{qmax\_mvar}；$P_{dc}$ 为双向推断对称界
        $\pm\max(10^{-3},\;|\,$\fld{pmax\_mw}$|,|\,$\fld{pmin\_mw}$|,|\,$
        \fld{qmax\_mvar}$|)/S_{\mathrm{base}}$（同文件:430--459）；三者初值均为
        0 向界内的截断（同文件:456--458）。
\end{itemize}
最后全体变量经 \srcpath{clamp\_in\_interior} 压入严格内点：边距
$\varepsilon=\min(\max(10^{-8},\,10^{-3}w),\;0.49w)$，$w$ 为该变量界宽
（同文件:205--223、461--463）；无法构造严格内点时直接失败返回
（同文件:3558--3562）。

\subsubsection*{目标函数}
目标为发电成本二次多项式之和，系数取自 canonical 发电机记录的
\fld{cost\_c2}/\fld{cost\_c1}/\fld{cost\_c0}（外部电网提升时同源拷贝），
\textbf{以 MW 有名值计价}
（\srcpath{src/optimal\_power\_flow/ac\_opf.cpp}:540--550，函数
\srcpath{evaluate\_objective}）：
\begin{align*}
f = \sum_{k=1}^{n_g}\left(c_{2,k}\,P_{g,k}^{2} + c_{1,k}\,P_{g,k} + c_{0,k}\right),
\qquad P_{g,k}\ \text{以 MW 计}
\end{align*}
内部变量 $P_g$ 为标幺值，故对 $x$ 的梯度与 Hessian 对角为
（同文件:552--572，函数 \srcpath{accumulate\_objective\_grad\_hdiag}）：
\begin{align*}
\frac{\partial f}{\partial P_{g,k}} &=
  2c_{2,k}S_{\mathrm{base}}^{2}\,P_{g,k} + c_{1,k}S_{\mathrm{base}} &
\frac{\partial^{2} f}{\partial P_{g,k}^{2}} &= 2c_{2,k}S_{\mathrm{base}}^{2}
\end{align*}
即目标对 $Q_g$、电压及换流器变量无直接二次项；无功与网损仅通过等式与不等式
间接进入目标。币种代码中未固定，单位记为``—''。

\subsubsection*{等式约束}
等式残差由共享潮流残差装配器与 OPF 自有项叠加而成
（\srcpath{src/optimal\_power\_flow/ac\_opf.cpp}:708--775，函数
\srcpath{assemble\_hybrid\_equalities}；潮流残差由
\srcpath{core::evaluate\_residual\_and\_jacobian} 在\textbf{剔除换流器}的
口径下计算，同文件:777--813）。

(1) 交流节点有功/无功平衡（$n_p+n_q$ 行，同文件:724--739）：
\begin{align*}
g^{P}_i &= P^{g}_i - P^{d}_i - P^{calc}_i(V,\theta) - \sum_{k\in\mathcal{C}(i)} P_{ac,k}
 & i &\in \text{非松弛母线}\\
g^{Q}_i &= Q^{g}_i - Q^{d}_i - Q^{calc}_i(V,\theta) - \sum_{k\in\mathcal{C}(i)} Q_{ac,k}
 & i &\in \text{PQ 母线}
\end{align*}
其中 $P^{calc}_i=V_i\sum_j V_j(G_{ij}\cos\theta_{ij}+B_{ij}\sin\theta_{ij})$、
$Q^{calc}_i=V_i\sum_j V_j(G_{ij}\sin\theta_{ij}-B_{ij}\cos\theta_{ij})$ 由 Ybus
直接求得（与 \srcpath{ac\_opf\_solver.hpp}:18--31 的文档式一致），
$\mathcal{C}(i)$ 为交流侧挂于母线 $i$ 的换流器集合。

(2) 直流节点平衡（$n_{dc}$ 行，同文件:741--755）：
\begin{align*}
g^{dc}_k &= \frac{P^{d}_{dc,k}}{S_{\mathrm{base}}}
  + V_{dc,k}\sum_{m} G^{dc}_{km}V_{dc,m} - \sum_{k'\in\mathcal{C}_{dc}(k)} P_{dc,k'}
\end{align*}
其中 $G^{dc}$ 为直流电导矩阵 \fld{gdc}，第一项为直流负荷，第二项为直流网络
注入 $V_k(G^{dc}V)_k$，末项扣除换流器直流口功率。

(3) 换流器交直流耦合（$n_c$ 行，同文件:757--772）：
\begin{align*}
g^{conv}_k &= P_{ac,k} + P_{dc,k} + P_{loss,k} = 0\\
P_{loss,k} &= a_k + b_k I_{ac,k} + c_k I_{ac,k}^{2}\\
I_{ac,k} &= \frac{\sqrt{\max(\varepsilon,\;P_{ac,k}^{2}+Q_{ac,k}^{2}+\varepsilon)}}
                    {\max(V_{m,k},\,10^{-6})}
\end{align*}
其中 $a_k=$~\fld{loss\_mw}$/S_{\mathrm{base}}$（固定损耗），
$b_k=$~\fld{loss\_percent}$/100$（比例损耗），$c_k=1-\eta$（\fld{eta} 效率，
导通损耗），数值软化常数 $\varepsilon=$~\fld{eps\_iac}=$10^{-6}$
（同文件:757--770）。符号约定为 $P_{ac}$ 交流口输出为正、
$P_{dc}=-(P_{ac}+P_{loss})$，即损耗由直流侧承担；该口径下\textbf{不计}
\fld{r\_conv\_ac\_pu} 交流侧导通损耗耦合（潮流模块才有该项），
\fld{converter\_model\_scope} 中 \fld{vsc\_ac\_conduction\_loss\_modelled}
保持 \fld{false}（同文件:3044--3050）。

\subsubsection*{非线性不等式（外点罚口径）}
紧致路径\textbf{不引入}不等式松弛与对应乘子，而是把支路热稳定界与换流器
容量圆写成 $h(x)\le 0$ 后以\textbf{外点二次罚}进入目标
（\srcpath{src/optimal\_power\_flow/ac\_opf.cpp}:853--966，函数
\srcpath{accumulate\_nonlinear\_limit\_terms}）：仅当 $h>0$ 时累加
\begin{align*}
\phi_{pen} &\mathrel{+}= \tfrac{1}{2}\rho\,h^{2}, &
\nabla\phi_{pen} &\mathrel{+}= \rho\,h\,\nabla h, &
\mathrm{diag}\,H_{pen} &\mathrel{+}= \rho\,(\nabla h)^{2}
\end{align*}
（同文件:865--885；Hessian 项为 Gauss--Newton 对角近似，舍去
$\rho h\nabla^{2}h$）。罚权重 $\rho$ 随迭代自适应，见
\ref{subsec:acopfnative:algorithm}~节。

(1) 支路视在功率平方界（两端各一条，同文件:887--952）。对在运且
\fld{rate\_a\_mva}$>0$ 的交流支路，以
$g_s=r/(r^{2}+x^{2})$、$b_s=-x/(r^{2}+x^{2})$、$b_c=b/2$、变比 $t$（0 回退 1）、
移相角 $\varphi=$~\fld{shift\_deg} 化弧度、
$\bar S^{2}=$(\fld{rate\_a\_mva}$/S_{\mathrm{base}})^{2}$
（\srcpath{build\_nonlinear\_limit\_workspace}，同文件:663--706）：
\begin{align*}
\theta &= \theta_i-\theta_j-\varphi\\
P_f &= \frac{g_s}{t^{2}}V_i^{2} - \frac{g_s\cos\theta+b_s\sin\theta}{t}V_iV_j\\
Q_f &= -\left(\frac{b_s}{t^{2}}+b_c\right)V_i^{2}
       - \frac{g_s\sin\theta-b_s\cos\theta}{t}V_iV_j\\
h_f &= P_f^{2}+Q_f^{2}-\bar S^{2}\le 0\\
P_t &= g_sV_j^{2} - \frac{g_s\cos\theta-b_s\sin\theta}{t}V_iV_j\\
Q_t &= -(b_s+b_c)V_j^{2} + \frac{g_s\sin\theta+b_s\cos\theta}{t}V_iV_j\\
h_t &= P_t^{2}+Q_t^{2}-\bar S^{2}\le 0
\end{align*}
（from 端：同文件:894--925；to 端：同文件:927--952；$\partial h/\partial x$
的四个分量在同文件:903--911、919--924、930--938、946--951 解析给出）。
串联阻抗为零（$r^{2}+x^{2}\le 0$）或未给 \fld{rate\_a\_mva} 的支路\textbf{跳过}
限值装配（同文件:667--678）。

(2) 换流器容量圆（每换流器一条，同文件:955--960）：
\begin{align*}
h^{conv}_k &= P_{ac,k}^{2}+Q_{ac,k}^{2}-\bar S^{2}_{conv,k}\le 0
\end{align*}
其中 $\bar S_{conv,k}$ 并非铭牌视在功率字段，而是从 P/Q 界推断的
$\max(|\,$\fld{pmax\_mw}$|,|\,$\fld{pmin\_mw}$|,|\,$\fld{qmax\_mvar}$|,|\,$
\fld{qmin\_mvar}$|)$，全部不可得时回退 \fld{kHugeBound}
（同文件:692--704）。

\subsubsection*{变量界的对数障碍}
盒式界经对数障碍进入目标（\srcpath{src/optimal\_power\_flow/ac\_opf.cpp}:
574--593，函数 \srcpath{accumulate\_log\_barrier}）：
\begin{align*}
\phi_{\mu} &= -\mu\sum_{i}\left[\ln(x_i-l_i)+\ln(u_i-x_i)\right]\\
\nabla\phi_{\mu,i} &= -\frac{\mu}{x_i-l_i}+\frac{\mu}{u_i-x_i} &
H_{\mu,ii} &= \frac{\mu}{(x_i-l_i)^{2}}+\frac{\mu}{(u_i-x_i)^{2}}
\end{align*}
任一分量触及或越出界时函数返回失败，主循环据以终止并报
``left the interior of variable bounds''（同文件:3641--3646）。

\subsection{原生内点算法全流程}\label{subsec:acopfnative:algorithm}

紧致路径的主循环为``外层障碍参数 $\mu$ 递减 + 内层障碍--罚 merit 牛顿法''
的双层结构（外层 \srcpath{src/optimal\_power\_flow/ac\_opf.cpp}:3609--4098，
内层 :3613--4081）。动态量初始化：
$\mu=\max($\fld{barrier\_mu0}$,$\fld{barrier\_mu\_min}$)$，
正则化 $\delta=\max($\fld{regularization}$,\,10^{-9})$，
罚权重 $\rho_{m}=\max($\fld{merit\_penalty}$,\,10)$，
目标权重 $w_f=0.05$（同文件:3596--3607）。

\subsubsection*{障碍--罚混合子问题}
每个内层迭代在当前 $(x,\mu,\rho_m,w_f)$ 下最小化
（同文件:3635--3678）：
\begin{align*}
F(x) = w_f f(x) + \phi_{\mu}(x) + \phi_{pen}(x;\rho_{nl})
       + \tfrac{1}{2}\rho_m\left(\lVert g(x)\rVert^{2}+\lVert h^{+}(x)\rVert^{2}\right)
\end{align*}
其中前四项进入梯度和 Hessian 对角（目标与障碍项精确，罚项取
Gauss--Newton 近似），末项仅以 merit 形式参与线搜索判据；非线性罚权重取
$\rho_{nl}=100\,\rho_m\cdot\mathrm{clamp}(10^{-2}/\max(\lVert g\rVert_{\infty},
10^{-12}),\,10^{-3},1)$，即可行域附近满罚、远离可行域时降罚以防梯度淹没
（同文件:3647--3650）。目标权重 $w_f$ 初值仅 0.05，在接近可行后逐步放大至
上限 1（同文件:3601、4078--4080）——这意味着迭代前中期求解的是一个
``可行性优先''的软化问题，判敛口径也随之针对障碍--罚子问题而非原始 KKT 系统。

\subsubsection*{KKT 牛顿系统组装与稀疏求解}
每次信赖尝试装配对称增广系统
（结构 \srcpath{build\_kkt\_workspace}，\srcpath{src/optimal\_power\_flow/ac\_opf.cpp}:1323--1399；
填值 \srcpath{update\_kkt\_values}，同文件:1401--1434）：
\begin{align*}
\begin{bmatrix} H+\delta I & J^{T}\\ J & -\delta I\end{bmatrix}
\begin{bmatrix}\Delta x\\ \Delta\lambda\end{bmatrix}
=
\begin{bmatrix}-\nabla F\\ -g\end{bmatrix}
\end{align*}
其中 $H$ 为目标+障碍+罚的对角 Hessian，$J=\partial g/\partial x$ 为等式雅可比，
右端 $\nabla F$ 不含 $J^{T}\lambda$ 项——对偶量 $\lambda$ 不跨迭代保持
（每次接受步后置零，同文件:4059），故该系统实质是``瞬时原始对偶''牛顿步；
解出的 $\Delta\lambda$ 仅随本步 merit 评估短暂留存，不作为对偶迭代量累积。
雅可比填值规则：潮流块取共享装配器输出的\textbf{负值}
（因 $g=$ 设定 $-$ 计算，注释见同文件:1236--1238）；发电机注入 $+1$、换流器
交流注入 $-1$（同文件:1241--1260）；直流块非对角元 $G^{dc}_{km}V_{dc,k}$、
对角元 $(G^{dc}V)_k+G^{dc}_{kk}V_{dc,k}$、$\partial g^{dc}/\partial P_{dc}=-1$
（同文件:1262--1283）；换流器耦合行
$\partial g/\partial P_{ac}=1+\partial P_{loss}/\partial P$、
$\partial g/\partial Q_{ac}=\partial P_{loss}/\partial Q$、
$\partial g/\partial P_{dc}=1$、$\partial g/\partial V_m=\partial P_{loss}/\partial V_m$，
其中 $\partial I/\partial P=P/(S_{ac}V_m)$ 等链式项解析给出
（同文件:1285--1320）。雅可比与 KKT 矩阵均采用``结构一次装配、非零元按索引
映射原地填值''的工作区复用设计（\srcpath{JWorkspace}:144--166、
\srcpath{KKTWorkspace}:168--176、\srcpath{build\_jacobian\_workspace}:968--1219）。
线性求解器由 \srcpath{core::make\_default\_sparse\_solver} 创建
（同文件:1397），后端按编译可用性取 KLU$\to$UMFPACK$\to$MKL
Pardiso$\to$SuperLU$\to$Eigen SparseLU 的优先级（
\srcpath{MIPSolvers/src/engine/kernel/linear\_algebra/linear\_solver.cpp}:230--242），
实际后端名记入 \fld{profiling.linear\_solver\_backend}（同文件:3551）。
稀疏模式分析只做一次并缓存，分解或回代失败即将本轮正则化乘 10 重试
（同文件:3694--3713、3835--3841）。

\subsubsection*{信赖尝试、步长截断与分数边界}
单个内层迭代最多做 \fld{kTrustAttempts}=6 次信赖尝试，每次乘以 10 递增正则化；
停滞计数 $\ge 3$ 且不可行度高于 $\max(10^{-3},\,10\,\varepsilon_{feas})$ 时跳过
信赖尝试直接进入恢复阶段（同文件:3686--3691）。牛顿步先做三层截断
（同文件:3715--3734）：
\begin{align*}
|\Delta x_i| &\le 0.25\,(u_i-l_i) && \text{逐变量}\\
\lVert\Delta x\rVert_{\infty} &\le 0.25 && \text{整体（\fld{kMaxDxInf}）}\\
|\Delta\lambda_r| &\le 10^{3} && \text{（\fld{kMaxDlambdaAbs}）}
\end{align*}
随后按分数边界规则取最大可行步长
（\srcpath{max\_feasible\_step}，同文件:1436--1454）：
\begin{align*}
\alpha_0 = \min\left(1,\;\min_{i:\Delta x_i>0}\frac{\tau(u_i-x_i)}{\Delta x_i},\;
                         \min_{i:\Delta x_i<0}\frac{\tau(x_i-l_i)}{-\Delta x_i}\right),
\qquad \tau=\mathrm{clamp}(\text{\fld{interior\_fraction}},\,0.5,\,0.9999)
\end{align*}
默认 $\tau=0.995$。

\subsubsection*{线搜索与接受判据}
在 $\alpha\leftarrow\alpha\cdot$\fld{step\_backoff}（默认 0.5）的回退链上最多试
\fld{max\_line\_search\_steps}（默认 20）个试验点，每个试验点要求严格内点、
重估等式与罚项（同文件:3742--3815）。接受判据为不可行度充分下降
\textbf{或}近可行时 merit 下降（同文件:3803--3812）：
\begin{align*}
\text{接受}\iff
\underbrace{\varphi^{trial}_{inf} < (1-5\times10^{-3}\alpha)\,\varphi^{cur}_{inf}}_{\text{不可行度下降}}
\;\lor\;
\big(\varphi^{cur}_{inf}\le\max(10^{-3},10\varepsilon_{feas})
 \;\land\; \varphi^{trial}_{merit}<\varphi^{cur}_{merit}\big)
\end{align*}
其中 $\varphi_{inf}=\sqrt{\lVert g\rVert^{2}+\lVert h^{+}\rVert^{2}}$。
若回退链上没有严格满足判据的点，启用``最优 $\alpha$ 兜底''：全程最佳试验点
在不可行度改善 $0.5\%$、或近可行且 merit 改善、或``近中性''（双指标恶化均
不超过 $0.5\%$）时仍可接受，近中性情形步长钳到 $5\times10^{-3}$
（同文件:3817--3833）。接受后正则化回收为
$\max($\fld{regularization}$,\delta_{trial}/2)$（同文件:3835--3837）。

\subsubsection*{恢复阶段与梯度回退}
信赖尝试全部失败时进入\textbf{恢复阶段}：最多 4 次尝试，Hessian 对角置零、
右端仅取 $-g$（纯可行性牛顿方向），步长上限收紧为
\fld{kMaxRestoreDxInf}=0.12 与逐变量 $0.10(u_i-l_i)$，接受判据放宽为
$(1-10^{-3}\alpha)$ 倍不可行度下降（同文件:3843--3958）。恢复再失败则取
最速可行性方向 $\Delta x=-J^{T}g$，上限 \fld{kMaxGradFallbackInf}=0.08，
判据 $(1-5\times10^{-4}\alpha)$ 或最佳点改善 $0.1\%$（同文件:3960--4035）。
三级机制均失败即终止，报 ``globalization failed (no acceptable
primal-dual/restoration step)''（同文件:4037--4041）。

\subsubsection*{步后更新与参数自适应}
接受步后：$x$ 压回内点（受约束变量边距比例 $10^{-3}$、自由变量 $10^{-8}$，
同文件:4043--4057）；$\lambda$ 置零（:4059）；按进展调整动态参数
（同文件:4061--4080）：
\begin{itemize}
  \item 恢复步：$\rho_m\leftarrow\min(10^{8},1.2\rho_m)$，
        $\delta\leftarrow\min(1,1.5\delta)$，$w_f\leftarrow\max(10^{-3},0.9w_f)$；
  \item 连续 4 步不可行度改善不足 $1\%$（停滞）：$\rho_m\times2$、$\delta\times5$、
        $w_f\times0.8$，停滞计数清零；
  \item 不可行度改善超 $20\%$：$\rho_m\leftarrow\max($\fld{merit\_penalty}$,
        0.95\rho_m)$（回落）；
  \item 不可行度进入 $10^{-2}$ 以内：$w_f\leftarrow\min(1,1.15w_f)$（目标权重爬坡）。
\end{itemize}

\subsubsection*{外层障碍参数更新与收敛判据}
内层收敛于当前 $\mu$（不可行度与平稳性同时达标，见下）或外层已``近可行''
（$\varphi_{inf}\le\max(10^{-4},\,50\varepsilon_{feas})$）时，外层把障碍参数
压向 $\mu\leftarrow\max(\mu\cdot$\fld{barrier\_mu\_reduction}$,\,$
\fld{barrier\_mu\_min}$)$（默认缩减率 0.2、下限 $10^{-8}$）；否则改为加重
$\rho_m\times1.5$、$\delta\times2$（同文件:4083--4098）。

收敛判据（同文件:3667--3674）：当
\begin{align*}
\max(\lVert g\rVert_{\infty},\;\max h^{+}) &\le \text{\fld{feasibility\_tol}}
&&\text{且} &
s_{dual} &\le \text{\fld{stationarity\_tol}}
\end{align*}
时认为当前 $\mu$ 的子问题收敛；若此时 $\mu\le1.01\,$\fld{barrier\_mu\_min}
则宣布整体收敛。平稳性度量 $s_{dual}$ 取 $\lVert\nabla F\rVert_{\infty}$；
当 $\varphi_{inf}\le10^{-3}$ 时改用其与最小二乘对偶估计
$\lVert\nabla F+J^{T}\lambda_{ls}\rVert_{\infty}$ 的较小者，$\lambda_{ls}$
由 $(JJ^{T}+10^{-8}I)\lambda_{ls}=-J\nabla F$ 经 \fld{SimplicialLDLT} 解得
（同文件:3655--3661、599--627）。

\subsubsection*{失败处理、诊断与回退}
终态汇总（\srcpath{finalize} 段，同文件:4100--4208）重估全部残差，按
$P/Q/DC/Conv$ 四族与 $\max h^{+}$ 分别报告，未收敛时 \fld{status} 携带各族
最大残差与最差平稳性变量的具名信息
（\srcpath{variable\_name}:1467--1499，状态串拼装:4196--4200）。失败原因枚举：
迭代耗尽（默认 ``maximum iterations reached''，:3597）、离开界内部（:3643）、
全局化失败（:4038）、初值非严格内点（:3558--3562）、维度退化或无在运发电机
（:3530--3539）。若 \fld{allow\_fallback} 且算例\textbf{不含} DC/VSC/LCC 组件，
未收敛时回退``经济调度 + AC 潮流''：先以 $\lambda$ 对分法做经济调度
（\srcpath{solve\_economic\_dispatch}:1501--1615，80 次对分、含截断退化的
最小平衡修正），再以牛顿潮流（80 次迭代、容差 $10^{-8}$）投影回交流可行域并
按母线均分 P/Q 到各机（\srcpath{try\_dispatch\_pf\_fallback}:1764--1867）；
混合算例则显式抑制该回退并写入提示（同文件:3266--3278、4189--4191）。
显式 \fld{EconomicDispatch} 后端直接使用同一快速路径，同样仅限纯 AC
（同文件:3501--3526）。

\subsubsection*{关键常数汇总}
\begin{center}
\footnotesize
\begin{tabular}{llll}
\toprule
常数 & 数值 & 出处（ac\_opf.cpp） & 含义\\
\midrule
\fld{kMinInteriorWidth} & $10^{-6}$ & :43 & 变量界最小宽度\\
\fld{kHugeBound} & $10^{4}$ & :44 & 自由变量数值界（标幺）\\
内点边距 & $\min(\max(10^{-8},10^{-3}w),0.49w)$ & :221 & 初值压入内点\\
$w_f$ 初值 & 0.05 & :3601 & 目标动态权重初值\\
$V_{dc}$ 界 & $[0.8,\,1.2]$ / DC\_V $\pm0.05$ & :418--423 & 直流电压硬编码界\\
\fld{eps\_iac} & $10^{-6}$ & :757 & 换流器电流数值软化\\
罚闸值 & $10^{-2}/\lVert g\rVert_\infty$ 钳于 $[10^{-3},1]$ & :3647 & 不等式罚自适应门\\
\fld{kTrustAttempts} & 6 & :3686 & 单迭代信赖尝试上限\\
\fld{kMaxDxInf} & 0.25 & :3687 & 原始步 $\infty$ 范数上限\\
\fld{kMaxDlambdaAbs} & $10^{3}$ & :3688 & 对偶步幅值上限\\
线搜索接受 & $(1-5\times10^{-3}\alpha)$ & :3804 & 不可行度充分下降系数\\
\fld{kRestorationAttempts} & 4 & :3844 & 恢复尝试上限\\
\fld{kMaxRestoreDxInf} & 0.12 & :3845 & 恢复步 $\infty$ 范数上限\\
\fld{kMaxGradFallbackInf} & 0.08 & :3963 & 梯度回退步上限\\
$JJ^{T}$ 正则 & $10^{-8}$ & :611 & 最小二乘对偶估计\\
ED 对分 & 80 次 / $10^{-8}$ & :1573--1581 & 经济调度 $\lambda$ 对分\\
PF 回退 & 80 次 / $10^{-8}$ & :1772--1773 & 回退潮流迭代与容差\\
\bottomrule
\end{tabular}
\end{center}

\subsection{Native parity IPM 核与求解编排}\label{subsec:acopfnative:kernel}

\subsubsection*{接口与数据结构}
作用于 Parity 全空间问题的自研原始对偶内点核声明于
\srcpath{include/hacdcpf/optimal\_power\_flow/native\_ipm\_solver.hpp}（命名空间
\fld{hacdcpf::opf::parity}，与 \fld{parity\_ipm\_solver.hpp} 共享结果类型）。
文档化的数学形式为（同文件:61--89）：
\begin{align*}
\min_x\; f(x)\quad \text{s.t.}\quad g(x)=0,\quad h(x)\le 0,\quad x_{min}\le x\le x_{max}
\end{align*}
不等式与盒式界引入松弛后，对摄动 KKT 系统
$\nabla f+J_g^{T}\lambda+J_h^{T}\mu+z^{+}-z^{-}=0$、$g=0$、$h+s=0$、$S\mu=\tau e$
取牛顿步，步长采用分数边界控制；返回对偶按 \fld{Problem::cidx} 的行序排列。
\fld{IPMOptions}（同文件:17--32）给出迭代上限 \fld{max\_iter}=400、三项容差
\fld{tol\_primal}/\fld{tol\_dual}/\fld{tol\_complementarity}=$10^{-6}$、正则化
$10^{-8}$、步长上限 \fld{alpha\_max}=0.95，以及可选的原始/等式对偶/不等式对偶/
松弛暖启动指针（仅在全部维数与正性检查通过时采用）。\fld{IPMResult}
（同文件:34--59）返回收敛标志、三项 KKT 残差及其初值、暖启动使用标记、
分解/回代/接受/拒绝/符号分析/后端升级/缩放重建计数、具体线性代数路径字符串
（如 \fld{dense\_lu}、\fld{sparse\_umfpack}、\fld{sparse\_klu}、
\fld{sparse\_eigen\_lu}）与原始-对偶解向量。核内部实现位于
\srcpath{src/optimal\_power\_flow/parity\_ipm.cpp}，细节见第~\ref{sec:parityipm}~章。
此外 \srcpath{homotopy\_tangent}（同文件:92--113）给出目标同伦 $P(t)=\min t\,f(x)$
在当前路径点的 Davidenko 切向：对增广系统右端取 $(-\nabla f/t,\,0,\,0)$ 复用
IPM 自身的 LDL$^{T}$ 机构一次分解回代，得 $(dx/dt,dz/dt,d\lambda/dt,d\mu/dt)$，
$t\approx0$、分解失败或非有限时返回 \fld{false}。

\subsubsection*{选项映射与暖启动编排}
\srcpath{solve\_with\_parity\_ipm}（\srcpath{src/optimal\_power\_flow/ac\_opf.cpp}:2486--3020）
把 \fld{ACOPFOptions} 翻译为核选项（同文件:2544--2551）：
\fld{max\_iter} 取内外层迭代上限之积；\fld{tol\_primal} 取
$\max($\fld{feasibility\_tol}$,10^{-10})$；\fld{tol\_dual} 取
$\max($\fld{stationarity\_tol}$,10^{-10})$；\fld{tol\_complementarity} 取
$\max($\fld{barrier\_mu\_min}$,10^{-10})$；\fld{regularization} 取
$\max($\fld{regularization}$,10^{-12})$；\fld{alpha\_max} 取
\fld{interior\_fraction} 钳于 $[0.5,0.9999]$。建模侧
\fld{load\_shedding} 仅当目标为 \fld{Economic} 时开启且 \fld{voll} 显式置 0，
四个 \fld{enforce\_*} 开关原样下传（同文件:2526--2540）。

暖启动有两条独立来源（同文件:2557--2602）：
\begin{itemize}
  \item \fld{ac\_pf\_warm\_start}（默认关）：先以混合潮流
        \srcpath{powerflow::solve\_hybrid}（开启换流器模式切换、PV-PQ 转换、
        自动摆节点选择）解出耦合可行点，再把电压、换流器功率、DC/DC 功率经
        \srcpath{build\_hybrid\_power\_flow\_warm\_start} 与
        \srcpath{recover\_power\_flow\_generator\_dispatch} 映射进 OPF 变量空间，
        其中发电机分派由潮流平衡残差在母线内按松弛优先重新分摊
        （同文件:2057--2265）；
  \item \fld{warm\_start}：把上一次兼容求解的 \fld{ACOPFResult} 逐变量族映射，
        优先采用不透明的 \fld{ipm\_primal\_state}（尺寸匹配且有限时），并以
        $10^{-8}$ 边距钳入界内；对偶与松弛状态仅在映射成功时一并传入
        （\srcpath{build\_parity\_primal\_warm\_start}:1879--2008）。
\end{itemize}

\subsubsection*{内层求解器选择与 Ipopt 回退}
内层引擎按 \fld{ParityInnerSolver}（枚举于同文件:1869--1873）与选项组合
（同文件:2604--2664）：
\begin{itemize}
  \item \fld{Auto} 且 \fld{ac\_pf\_warm\_start} 且非同伦：先用 Ipopt（若编译入），
        其终止于 acceptable 带而未严格收敛时，把最末原始点连同对偶/松弛交给
        自研核精化；Ipopt 不可用时直接自研核；
  \item \fld{Auto} 其余情形：先自研核，未收敛且 \fld{allow\_fallback} 时回退 Ipopt；
  \item \fld{NativeIPM}：仅自研核；\fld{Ipopt}：仅 Ipopt，未编译时回退自研核。
\end{itemize}
Ipopt 适配（\srcpath{solve\_parity\_with\_ipopt}:2384--2484，仅在
\srcpath{HACDCPF\_HAVE\_IPOPT} 下编译）把 Parity 问题的目标/梯度/等式/不等式/
雅可比/拉格朗日 Hessian 回调一一挂到 \fld{engine::NLPModel}
（\srcpath{build\_parity\_nlp\_input}:2275--2372），容差映射为：主容差取三项
最小、acceptable 带取三项最大、acceptable 互补下限 $\max(10^{-2},100\,
$\fld{tol\_complementarity}$)$（同文件:2321--2338）。结果接受口径除
\fld{success} 外还接受``acceptable residuals''：原始残差
$\le$\fld{tol\_primal}、对偶 $\le1000\,$\fld{tol\_dual}、互补
$\le\max(10^{-2},1000\,$\fld{tol\_complementarity}$)$，状态串保留 Ipopt 原始终止码
（同文件:2412--2428、2472--2475）。对偶重建时非线性不等式与盒式界乘子以
$2\times10^{-12}$ 为下限拼回 \fld{mu}/\fld{z}（同文件:2435--2471）。

\subsubsection*{目标同伦延续}
\fld{objective\_homotopy} 开启且内层非 Ipopt 时，进入两阶段延续
（同文件:3352--3482；理论引用 docs/numerical\_methods.md §12--13）：沿
$P(t)=\min\,t\,f$ 从 $t=0$（纯可行性）到 $t=1$（全成本），成本系数与 RPO 目标
权重同步乘 $t$（同文件:3371--3385）；最多 120 步，步长初值
\fld{homotopy\_dt0}（默认 0.10），快收敛（$\le20$ 次迭代）步长加倍至上限 0.5、
失败减半，$dt<0.005$ 判定进入最小违反域并停止（同文件:3367--3368、3395--3398、
3437--3443）。相邻步之间用 Davidenko 切向做预测
$w_{pred}=w+h\,w'$，预测步长受参数度量信赖半径 $\kappa=0.05$ 约束
$h=\min(dt,\kappa/\lVert dx/dt\rVert_\infty)$，松弛与不等式对偶以 $10^{-8}$ 为底
（同文件:3399--3436）。$t=1$ 未达且 \fld{allow\_fallback} 时，从全新耦合潮流点
对全目标做最后一次独立求解，成功则标记
``converged (full objective after incomplete objective homotopy)''；否则诚实报告
``not converged: objective homotopy stopped at t=...''（同文件:3451--3478）。
显式 Ipopt 请求与同伦互斥：同伦标志被清除后做单次独立求解（同文件:3483--3488）。

\subsubsection*{结果提取、LMP 与诊断量}
求解返回后（同文件:2666--3019）：电压、发电、甩负荷、换流器、DER、储能、
DC/DC、柔性负荷、能量路由器端口按块从原始向量提取并乘回 MW/MVAr
（:2716--2852）；在运 LCC 在最优电压状态下重估准稳态工作点
（$\alpha$/$\gamma$/$U_{d0}$/$I_d$/越限审计），触发角或熄弧角越窗时写入
\fld{model\_limitations}（:2854--2926）；LMP 从等式对偶提取
$\lambda^{P}_i\cdot$\fld{scale\_p}$/S_{\mathrm{base}}$ 与
$\lambda^{Q}_i\cdot$\fld{scale\_q}$/S_{\mathrm{base}}$，仅当对偶向量完整覆盖
平衡行时置 \fld{lmp\_valid=true}，否则写明原因（Ipopt 适配器未返回约束乘子或
对偶不完整，:2930--2952）；终态按平衡族报告 pu 残差并以
$\max(\lVert g\rVert_\infty,\max h^{+},\text{界违反})$ 汇总
\fld{max\_constraint\_violation}（:2954--2996）；原始/对偶/松弛终态存入
\fld{ipm\_*\_state} 供续解（:2684--2695）；同伦切向存入 \fld{ipm\_tangent\_*}
（:2697--2714）。最后按 \fld{bus\_merge\_map} 反投影到用户母线序并按投影证书
裁剪内部直流母线（:3011--3017；\srcpath{unproject\_per\_bus\_ac\_opf\_result}
与 \srcpath{trim\_internal\_dc\_bus\_results}，同文件:46--77）。计数口径：
\fld{iterations} 报告为 IPM 迭代数 $+1$、\fld{outer\_iterations} 恒为 1
（:2667--2668）。

\subsection{公共接口与门面}\label{subsec:acopfnative:api}

本引擎对外的公开 API 均为\textbf{自由函数}（当前代码检出处不存在名为
\fld{NativeIpmSolver} 或 \fld{AcOpfSolver} 的类）：
\begin{itemize}
  \item \texttt{opf::solve\_ac\_opf(sys, opt)}：主入口
        （\srcpath{include/hacdcpf/optimal\_power\_flow/ac\_opf\_solver.hpp}:42--43）；
  \item \texttt{opf::compute\_jacobian\_diagnostics(sys, opt)}：雅可比诊断声明
        （同文件:54--55）。\textbf{注意}：该函数在当前检出处\textbf{只有声明、
        没有定义}；\srcpath{ac\_opf.cpp}:4211 实现的是
        \srcpath{check\_ac\_opf\_jacobian\_fd}，且无任何调用方，两者均未接入
        构建产物（详见 \ref{subsec:acopfnative:verification}~节）；
  \item \texttt{parity::solve\_primal\_dual\_ipm(prob, opt)} 与
        \srcpath{parity::homotopy\_tangent(...)}：Native IPM 核与同伦切向
        （\srcpath{native\_ipm\_solver.hpp}:90--113）；
  \item \fld{hacdcpf::solve\_ac\_opf} 全局门面
        （\srcpath{src/api/hacdcpf.cpp}:1865--1890）：调用 \fld{opf::solve\_ac\_opf}
        后，对收敛且含开关/断路器的算例补算开关类设备端子潮流
        （\fld{ac\_switch\_flows}/\fld{ac\_circuit\_breaker\_flows}）；
  \item \fld{IOPFSolver}/\fld{DefaultOPFSolver} 抽象
        （\srcpath{include/hacdcpf/optimal\_power\_flow/opf\_solver\_interface.hpp}:24--42）：
        把 \fld{OPFProblem} 路由到 \fld{hacdcpf::solve\_ac\_opf} 或
        \fld{solve\_dc\_opf}，\fld{sys} 为空或 \fld{is\_ac}/\fld{is\_dc} 不唯一时
        抛 \fld{std::invalid\_argument}。
\end{itemize}

\subsection{选项字段表}\label{subsec:acopfnative:options}

下表给出本章引擎直接消费的 \fld{ACOPFOptions} 字段（定义于
\srcpath{include/hacdcpf/optimal\_power\_flow/opf\_options.hpp}:49--123；
公共结果字段 \fld{converged}/\fld{objective}/\fld{status} 等的通用口径见
第~\ref{sec:overview}~章）。``$\dagger$''标记仅 Parity 路径消费的字段。

\begin{paramtable}
\fld{ac\_solver\_backend} & — & — & Auto/\allowbreak ParityIPM/\allowbreak Ipopt/\allowbreak EconomicDispatch & Auto & 后端选择；翻译规则见 \ref{subsec:acopfnative:role}~节\\
\fld{objective} & — & — & Economic/\allowbreak VoltageDeviation/\allowbreak ActiveLoss/\allowbreak VoltageDeviation\allowbreak AndLoss & Economic & 目标类别；非经济目标经 Parity 建模生效$\dagger$\\
\fld{voltage\_target\_pu} & $V^{ref}$ & pu & 0.95--1.05（经验值） & 1.0 & 电压偏差目标$\dagger$\\
\fld{voltage\_deviation\_weight} & $w_V$ & — & $>0$ & 1.0 & 电压偏差目标权重$\dagger$\\
\fld{active\_loss\_weight} & $w_{loss}$ & — & $>0$ & 1.0 & 网损目标权重$\dagger$\\
\fld{max\_inner\_iterations} & — & — & $\ge1$ & 80 & 内层迭代上限；$\le0$ 时重置 80\\
\fld{max\_outer\_iterations} & — & — & $\ge1$ & 8 & 外层 $\mu$ 迭代上限；$\le0$ 时重置 8\\
\fld{max\_line\_search\_steps} & — & — & $\ge1$ & 20 & 线搜索回退上限\\
\fld{feasibility\_tol} & $\varepsilon_{feas}$ & pu & $10^{-8}$--$10^{-4}$ & $10^{-6}$ & 等式/不等式违反判敛容差\\
\fld{stationarity\_tol} & $\varepsilon_{stat}$ & — & $10^{-8}$--$10^{-4}$ & $10^{-6}$ & 平稳性判敛容差\\
\fld{barrier\_mu0} & $\mu_0$ & — & $10^{-4}$--1 & $10^{-2}$ & 初始障碍参数\\
\fld{barrier\_mu\_reduction} & $\sigma_\mu$ & — & 0.1--0.5 & 0.2 & 外层 $\mu$ 缩减率\\
\fld{barrier\_mu\_min} & $\mu_{min}$ & — & $10^{-10}$--$10^{-6}$ & $10^{-8}$ & 障碍参数下限（收敛门槛）\\
\fld{merit\_penalty} & $\rho_m$ & — & 1--$10^{3}$ & 10.0 & merit 等式罚权重初值\\
\fld{regularization} & $\delta$ & — & $10^{-9}$--$10^{-4}$ & $10^{-6}$ & KKT 正则化初值\\
\fld{step\_backoff} & $\beta$ & — & 0.3--0.7 & 0.5 & 线搜索回退因子\\
\fld{interior\_fraction} & $\tau$ & — & 0.9--0.9999 & 0.995 & 分数边界比例\\
\fld{ac\_eval\_threads} & — & — & $\ge1$ & 1 & 交流残差并行线程数\\
\fld{enable\_primal\_dual} & — & — & true/false & false & legacy 标志：启用内点路径（Auto 下混合算例自动置位）\\
\fld{use\_parity\_ipm} & — & — & true/false & false & legacy 标志：改用 Parity 全空间建模\\
\fld{allow\_fallback} & — & — & true/false & true & 允许 Ipopt/ED+PF 回退与同伦末段挽救\\
\fld{verbose} & — & — & true/false & false & 迭代日志\\
\fld{enable\_phase\_one} & — & — & true/false & true & Parity Native IPM 的有界 primal/dual warm-start 构造$\dagger$\\
\fld{phase\_one\_time\_limit\_ms} & — & ms & $>0$ 协作截止；$\le0$ 不限 & 5000 & Phase I 墙钟预算；单次稀疏分解不可中断$\dagger$\\
\fld{phase\_one\_max\_iterations} & — & — & $\ge0$ & 12 & Phase I Newton 迭代硬上限$\dagger$\\
\fld{phase\_one\_max\_factorizations} & — & — & $\ge0$ & 14 & primal/dual 共用稀疏分解硬上限$\dagger$\\
\fld{phase\_one\_max\_backtracks} & — & — & $\ge0$ & 12 & 每次 Phase I 步的回溯上限$\dagger$\\
\fld{phase\_one\_dispatch\_dual\_predictor} & — & — & true/false & false & 实验性零分解连通分量 dispatch-price dual 候选$\dagger$\\
\fld{phase\_one\_dispatch\_dual\_min\_improvement} & — & — & 0--1 & 0.05 & raw 与 normalized dual residual 必须同时达到的改善比例$\dagger$\\
\fld{warm\_start} & — & — & 指针/null & null & 上次兼容结果的非拥有指针，求解期间须存活\\
\fld{ac\_pf\_warm\_start} & — & — & true/false & false & 先做混合潮流并映射为暖启动（默认关，注释称良态算例物理初值更快）\\
\fld{compute\_homotopy\_tangent} & — & — & true/false & false & 返回点计算 Davidenko 切向$\dagger$\\
\fld{homotopy\_t} & $t$ & — & $(0,1]$ & 1.0 & 所解问题的成本标度（切向右端 $\nabla f/t$）$\dagger$\\
\fld{objective\_homotopy} & — & — & true/false & false & 两阶段目标同伦延续；与 Ipopt 内层互斥\\
\fld{homotopy\_dt0} & $dt_0$ & — & 0.01--0.5 & 0.10 & 同伦初始步长；$\le0$ 回退 0.10\\
\fld{enforce\_branch\_limits} & — & — & true/false & true & 支路限值族开关$\dagger$\\
\fld{enforce\_converter\_capacity} & — & — & true/false & true & 换流器容量圆开关$\dagger$\\
\fld{enforce\_converter\_current\_limits} & — & — & true/false & true & 换流器电流限值开关$\dagger$\\
\fld{enforce\_converter\_modulation\_limits} & — & — & true/false & true & 调制度/占空比窗口开关$\dagger$\\
\end{paramtable}

\subsection{结果字段表}\label{subsec:acopfnative:results}

\fld{ACOPFResult}（\srcpath{include/hacdcpf/optimal\_power\_flow/opf\_result.hpp}:72--169）
字段众多；下表仅列与本章引擎直接相关者，派生设备族（DER/储能/DC-DC/柔性负荷/
能量路由器，仅 Parity 路径填充）与通用字段的完整口径见第~\ref{sec:overview}~章。

\begin{paramtable}
\fld{vm} / \fld{va} & $V,\theta$ & pu, rad & 约 0.9--1.1 / $\pm\pi$ & 1.0 / 0 & 交流母线电压（反投影回用户母线序）\\
\fld{vdc} & $V_{dc}$ & pu & 0.8--1.2 & 空 & 直流母线电压（按投影证书裁剪）\\
\fld{pg\_mw} / \fld{qg\_mvar} & $P_g,Q_g$ & MW, MVAr & 机界内 & 输入值 & 发电机调度（已剔除合成外部电网机）\\
\fld{external\_grid\_p\_mw} / \fld{external\_grid\_q\_mvar} & $P_{eg},Q_{eg}$ & MW, MVAr & $\pm$合成界 & 0 & 外部电网分派（由合成机拆回）\\
\fld{dpd\_mw} / \fld{dqd\_mvar} & $\Delta P_d,\Delta Q_d$ & MW, MVAr & $\ge0$ & 空 & 逐母线甩负荷（仅 Parity 路径）\\
\fld{pac\_mw} / \fld{qac\_mvar} & $P_{ac},Q_{ac}$ & MW, MVAr & 换流器界内 & 空 & VSC 交流口功率（仅 Parity 路径填充）\\
\fld{lcc\_transfers} & — & — & — & 空 & 最优电压点重估的 LCC 准稳态工作点（仅 Parity 路径）\\
\fld{lmp\_p} / \fld{lmp\_q} & $\lambda^P,\lambda^Q$ & — & — & 空 & 节点价格，$\lambda\cdot$scale$/S_{\mathrm{base}}$；币种未定\\
\fld{lmp\_valid} / \fld{lmp\_validity\_reason} & — & — & true/false & false/空 & LMP 是否来自认证等式乘子及原因串\\
\fld{ipm\_primal\_state} 等四元 & — & — & — & 空 & 不透明续解状态（原始/等式对偶/不等式对偶/松弛），配合 \fld{warm\_start} 使用\\
\fld{ipm\_tangent\_primal} 等四元 & $dw/dt$ & — & — & 空 & 同伦切向（\fld{compute\_homotopy\_tangent} 时填充）\\
\fld{gen\_map} & — & — & — & — & 发电机行到 \fld{index} 的稳定归因映射\\
\fld{iterations} / \fld{outer\_iterations} & — & — & $\ge0$ & 0 & 迭代计数（Parity 路径为 IPM 数+1 与恒 1）\\
\fld{max\_constraint\_violation} & $\varphi_{inf}$ & pu & — & 0 & 终态最大约束违反（等式、$h^+$、界三者取大）\\
\fld{max\_stationarity} & $s_{dual}$ & — & — & 0 & 终态平稳性度量（障碍-罚口径，近可行时取最小二乘估计）\\
\fld{solver\_path} & — & — & NativeAC/\allowbreak ParityIPM & Unknown & 实际求解路径标记\\
\fld{profiling.linear\_solver\_backend} & — & — & 后端名字符串 & 空 & KKT 线性代数路径\\
\fld{profiling.\allowbreak factorization\_calls} 等计数 & — & — & $\ge0$ & 0 & 分解/回代/接受/拒绝/分析/升级/缩放重建计数\\
\fld{profiling.phase\_one\_*} & — & pu/ms/计数 & — & inf/0/not-run & Phase I 前后违反量、dual-fit、预算消耗、终止/线性后端及 Phase II 接受状态\\
\fld{profiling.dispatch\_dual\_predictor\_*} & — & ms/残差/布尔 & — & inf/0/not-attempted & 零分解 dual predictor 的尝试、接受、成本及 exact residual 审计\\
\fld{profiling.max\_*\_residual\_pu} & — & pu & — & 0 & 分族终态残差（P/Q/DC/Conv/其他/非线性不等式）\\
\fld{profiling.final\_barrier\_mu} & $\mu$ & — & — & 0 & 终态障碍参数（Parity 路径记录的是互补度）\\
\fld{infeasibility\_hints} & — & — & — & 空 & 不可行/不收敛的可能原因\\
\fld{model\_limitations} & — & — & — & 空 & 诚实口径声明（非凸局部解、储能单时段等）\\
\fld{converter\_model\_scope} & — & — & scope 标签 & 空 & 换流器建模覆盖声明（NativeAC 为 \texttt{ac-opf-\allowbreak native:\allowbreak vsc-free-\allowbreak pq+\allowbreak capacity-\allowbreak circle}）\\
\fld{audit} & — & — & — & 空 & 求解后独立可行性审计（由 \fld{verify\_opf\_result} 填充，见第~\ref{sec:verification}~章）\\
\fld{ac\_switch\_flows} / \fld{ac\_circuit\_breaker\_flows} & — & MW 等 & — & 空 & 开关/断路器端子潮流（全局门面补算）\\
\end{paramtable}

\subsection{验证与校核}\label{subsec:acopfnative:verification}

\begin{itemize}
  \item \textbf{解析雅可比的有限差分自检}：
        \srcpath{check\_ac\_opf\_jacobian\_fd}（\srcpath{src/optimal\_power\_flow/ac\_opf.cpp}
        第 4211--4409 行）对全部或均匀抽样的变量列做中心差分
        $h=\min($\fld{fd\_eps}$\cdot\max(1,|x|),\;0.49\min(x-l,u-x))$，
        \fld{fd\_eps} 缺省 $10^{-6}$，逐元比较解析列与差分列，通过阈值为
        最大绝对误差 $\le5\times10^{-4}$ 或最大相对误差 $\le5\times10^{-3}$
        （同文件:4319--4407）。\textbf{如实说明}：该函数在当前检出处没有任何
        调用方；公共头声明的 \srcpath{compute\_jacobian\_diagnostics}
        （\srcpath{ac\_opf\_solver.hpp}:54--55）在 \srcpath{src/} 内无定义，
        两者名字不一致且均未接入默认测试目标，属已登记的存疑点。
  \item \textbf{后端选择与混合算例收敛}：\srcpath{tests/test\_opf\_solver\_backends.cpp}
        覆盖内部混合算例 \fld{case300\_acdc}/\fld{case2000\_acdc} 的 ParityIPM
        收敛（:86 起）、Ipopt 环境变量门控，以及 \fld{case2869pegase} 大算例的
        稀疏 KKT 路径收敛（:1154 起）。
  \item \textbf{三方法交叉验证}：\srcpath{tests/test\_acopf\_dcopf\_crossval.cpp}
        在 MATPOWER 小算例（case5/9/14/30/39/57）上对 AC OPF、DC OPF 与潮流
        三者做调度一致性互查（:126、:508、:637 调用 \fld{solve\_ac\_opf}）。
  \item \textbf{回归与特殊场景}：\srcpath{tests/test\_hacdcpf.cpp}:390 的
        IEEE14 AC/DC 混合 OPF；\srcpath{tests/test\_nighttime\_opf.cpp}:58 起的
        夜间无光照 OPF 回归；\srcpath{tests/test\_graph\_roundtrip.cpp}:249 的
        图聚合往返后 OPF 一致性。
  \item \textbf{手算校核数据}：\srcpath{external\_data/opf\_hand\_cases/} 提供
        单母线 merit-order、两母线拥塞、甩负荷、二次成本线性化四个手算可复算
        算例（见其 README）；当前检出处未见直接消费该目录的 C++ 测试目标，
        如实标注为数据资产。
  \item \textbf{大网实测声明（代码注释口径）}：\fld{ac\_pf\_warm\_start} 与
        \fld{objective\_homotopy} 的字段注释记载 PEGASE 大网实测——
        case13659pegase 经暖启动以 189 次迭代收敛于参考最优、同伦延续认证
        2869 节点精确解与 9241 节点 $-0.013\%$ 偏差
        （\srcpath{opf\_options.hpp}:88--114）。这些数字为注释记载、非本章实测复核。
  \item \textbf{万节点结构化启动实测}：Release/KLU、一次预热加三次计时的固定协议下，
        case9241pegase 从普通 AC-PF 启动的 $12.657\,\mathrm{s}$/77 次迭代/76 次
        Phase-II 分解，降至 compact DCOPF dispatch seed 的
        $11.088\,\mathrm{s}$/56/55；端到端中位数下降 $12.4\%$，分解下降
        $27.6\%$。两者目标分别为 315911.571480 与 315911.571406，终态
        primal/dual 均低于 $5\times10^{-7}$。case13659pegase 的 baseline dual
        为 1.970，预筛选零成本跳过 DC，保持 138 次迭代、137 次分解、目标
        386116.837520 及 $2.05\times10^{-7}/5.89\times10^{-9}$ 的最终证书。
        该路径按真实 Parity 初始指标选择 dispatch 点；9241 的 dual
        $838.166\to236.753$，但最大 primal metric
        $1.601\times10^{-3}\to1.506\times10^{-1}$，故它是 dual-dispatch
        warm start 而非 central-state 证书，最终 KKT 仍由 Phase II 独立认证。
        DC Phase I 的 $2000\,\mathrm{ms}$ 默认墙钟预算覆盖投影、formulation、
        结构初值、symbolic 与 numeric 工作；本例耗时 $1638\,\mathrm{ms}$、一次
        DC symbolic analyze，且 baseline/candidate 只共享构建一次的 Parity
        formulation。稀疏分解不可中断，故允许一个在途分解造成的可诊断超时。
  \item \textbf{两万节点级边界实测}：25,000-bus ACTIVSg25k（4,834 台发电机、
        32,230 条交流支路）在同一 Release/KLU 主机上按一次预热加三次计时协议，
        普通 AC-PF 启动的样本为 $183.967/182.383/183.808\,\mathrm{s}$，中位数
        $183.808\,\mathrm{s}$；三次均以 66 次迭代、65 次 Phase-II 分解收敛，
        目标为 6252070.20694，primal/dual 为
        $3.239\times10^{-7}/1.004\times10^{-7}$。默认 $2\,\mathrm{s}$ DC Phase I
        的中位数为 $186.588\,\mathrm{s}$（慢 $1.51\%$）：一次 symbolic、两次
        迭代后以 $2064\,\mathrm{ms}$ 超时（overshoot $64\,\mathrm{ms}$），未形成
        合格候选并精确回退，故 Phase II 工作量与最终 KKT 不变。
        将 DC 预算扩大到 $30\,\mathrm{s}$ 虽把 DC 残差降至
        $8.767\times10^{-4}$，但真实 Parity primal 从
        $3.929\times10^{-4}$ 恶化到 $9.786\times10^{-2}$，dual 仅改善
        $4.15\%$，仍被固定 admission gate 拒绝，总时长反而增至
        $215.012\,\mathrm{s}$。因此当前实现已验证能求解该 25k ACOPF，但不能
        声称结构化 DC Phase I 在该算例上加速；继续提高辅助 DC-IPM 精度不是
        合理方向。
  \item \textbf{零分解 dual predictor 负结果}：实验模式按交流导通分量从未贴界
        机组 stationarity 中位数构造共同 $P/Q$ 平衡乘子，不增加 symbolic/numeric
        分解，并要求 raw 与 normalized full dual residual 均改善至少 $5\%$。
        ACTIVSg25k 单次 Release/KLU 探针仅以 $1.978\,\mathrm{ms}$ 将二者从
        $47278.840/468.10733$ 改为 $47276.850/468.08762$（约 $0.0042\%$），
        因而拒绝；总时长 $184.344\,\mathrm{s}$，Phase II 仍为 66 次迭代、
        65 次分解。该选项默认关闭，不能作为 25k 加速声明。
  \item \textbf{25k prepared-session 重复求解}：同一 ACTIVSg25k 在一个
        \srcpath{PreparedACOPFSession} 内连续求解两次，首轮 AC-PF 启动耗时
        $184.645\,\mathrm{s}$、66 次迭代、65 次 numeric 分解；第二轮复用
        formulation/maps、symbolic ordering 与完整 continuation state，耗时
        $8.666\,\mathrm{s}$、4 次迭代、3 次新 numeric 分解、0 次 symbolic
        analyze，墙钟降低 $95.31\%$。第二轮目标 6252069.93711，primal/dual
        为 $3.389\times10^{-8}/2.092\times10^{-8}$；峰值 RSS 为
        $4709.6\,\mathrm{MiB}$。因此 25k 重复求解加速已验证，但内存效率仍是
        更大批处理的边界；该结果不表示复用了旧 numeric factors。
\end{itemize}

\subsection{边界与局限}\label{subsec:acopfnative:limitations}

以下条目均按代码实际行为（含其自述注释与 \fld{model\_limitations} 文案）列出：
\begin{enumerate}
  \item \textbf{非凸局部解}：AC OPF 为非凸非线性规划，收敛仅认证局部 KKT 点而非
        全局最优——该语句作为 \fld{model\_limitations} 固定写入每次返回
        （\srcpath{src/optimal\_power\_flow/ac\_opf.cpp}:3205--3206）。
  \item \textbf{legacy 定位}：紧致 NativeAC 路径自述为旧式``外部障碍/罚循环''，
        结果中建议生产环境改用 ParityIPM（同文件:3219--3222）。
  \item \textbf{不等式非硬约束}：支路热稳定界与容量圆仅经外点二次罚进入目标，
        终解允许受罚权重控制的残余违反；罚 Hessian 为 Gauss--Newton 对角近似，
        且罚权重在远离可行域时按门值降权（同文件:853--966、3647--3650）。
  \item \textbf{无持续对偶迭代}：$\lambda$ 每次接受步后置零（同文件:4059），
        平稳性判据针对障碍--罚软化子问题，近可行时才切换到最小二乘对偶估计；
        目标权重 $w_f$ 初值 0.05、随可行性爬坡，前中期迭代并非在解原始成本问题
        （同文件:3601、3655--3661、4066--4080）。
  \item \textbf{换流器控制模式无关}：紧致路径把 VSC 处理为自由 $P_{ac}/Q_{ac}/P_{dc}$
        盒式变量，不锚定 VDC\_Q 直流电压、不实施电流/调制度窗口，
        \fld{vsc\_vdc\_control\_modelled=false}，scope 标签固定为
        \fld{ac-opf-native:vsc-free-pq+capacity-circle}（同文件:3041--3050）；
        亦不计 \fld{r\_conv\_ac\_pu} 交流侧导通损耗。
  \item \textbf{推断界与硬编码界}：换流器容量圆半径与 $P_{dc}$ 界由 P/Q 界推断
        而非铭牌视在功率（同文件:443--447、692--704）；$V_{dc}$ 界硬编码
        $[0.8,1.2]$、DC\_V 型收缩为设定值 $\pm0.05$（:415--428）；$\theta$ 界
        硬编码 $\pm\pi$（:378）；无显式 SLACK 时取第 0 号母线（:178--185）。
  \item \textbf{LCC 不进入紧致路径}：无 $V_{dc}$ 耦合列，凡在运 LCC 一律强制
        Parity/Ipopt 路径（同文件:3343--3350）；Parity 路径中 LCC 的
        $P/I/\alpha/\gamma$ 指令与换流变分接头为固定输入而非调度变量，额定电流
        饱和纳入模型而 $\alpha/\gamma$ 限值仅为求解后审计，换相重叠迭代、
        分接头控制与有功损耗未建模，等式二阶导用解析雅可比的局部有限差分且
        额定电流转折不可微（同文件:2514--2523、2854--2926）。
  \item \textbf{设备覆盖边界}：储能按单时段功率注入优化、不含跨时段 SOC 动态；
        能量路由器端口平衡按无损处理（同文件:3207--3214）；支路限值仅取
        \fld{rate\_a\_mva}（:667）；甩负荷仅 Economic 目标开启且 \fld{voll}
        置 0，其价格口径由 Parity 建模层决定（:2527--2531，见
        第~\ref{sec:parityform}~章）。
  \item \textbf{外部电网合成界}：短路容量并非进口限值，仅作有限尺度；该近似已在
        代码注释中自述（同文件:3173--3189）。
  \item \textbf{回退口径}：ED+PF 回退对混合算例显式抑制（同文件:3266--3278、
        4189--4191）；回退成功时 \fld{max\_stationarity} 置 0、目标按钳制后出力
        重算，属可行性投影而非优化解（:1832--1845）。
  \item \textbf{门控依赖}：Ipopt 路径仅在 \srcpath{HACDCPF\_HAVE\_IPOPT} 编译时
        存在，且可被 \srcpath{HACDCPF\_DISABLE\_IPOPT\_OPF} 运行时关闭
        （同文件:33--35、2388--2396）；未编译时 \fld{Auto} 的全部回退自动失效。
  \item \textbf{已知不一致}：\srcpath{compute\_jacobian\_diagnostics} 有声明无定义、
        \srcpath{check\_ac\_opf\_jacobian\_fd} 有定义无调用方；
        \srcpath{rebalance\_pg\_within\_bus} 标记 \fld{[[maybe\_unused]]} 为死代码
        （同文件:1617）；回退用 \srcpath{compute\_ac\_injections} 为稠密 Ybus
        双重循环，仅小算例回退可接受（:1738--1762）。
\end{enumerate}
